% Syntra LaTeX Template
% For Arabic support, use XeLaTeX with:
%   \usepackage{fontspec}
%   \usepackage{polyglossia}
%   \setdefaultlanguage{english}
%   \setotherlanguage{arabic}
%   \newfontfamily\arabicfont[Script=Arabic]{Amiri}

\documentclass[12pt,a4paper]{article}

% Font and encoding
\usepackage[utf8]{inputenc}
\usepackage[T1]{fontenc}

% Math packages
\usepackage{amsmath,amssymb}

% Table support
\usepackage{longtable,booktabs}

% Graphics and links
\usepackage{graphicx}
\usepackage{hyperref}
\usepackage{geometry}

% Page geometry
\geometry{
    left=2.5cm,
    right=2.5cm,
    top=3cm,
    bottom=3cm
}

% Fix pandoc tightlist issue
\providecommand{\tightlist}{\setlength{\itemsep}{0pt}\setlength{\parskip}{0pt}}

% Hyperref settings
\hypersetup{
    colorlinks=true,
    linkcolor=blue,
    filecolor=magenta,
    urlcolor=cyan,
    pdftitle={Syntra Research Documentation},
    pdfauthor={Syntra-Env},
    pdfpagemode=FullScreen
}

\begin{document}

\section{Syntra}\label{syntra}

\emph{Towards Free Agency for All.} Syntra explores how to make personal
AI agency universally accessible --- lightweight agents that run on edge
devices, protect users from psychological manipulation, and provide
proactive guidance without being intrusive.

\subsection{Core Ideas}\label{core-ideas}

\begin{itemize}
\tightlist
\item
  \textbf{Edge-first agents} --- fast, local models that run on consumer
  hardware without cloud dependency.
\item
  \textbf{Manipulation defense} --- agents that identify and shield
  users from dark patterns, deceptive design, and persuasion techniques.
\item
  \textbf{Proactive but respectful} --- useful guidance that surfaces at
  the right time without becoming noise.
\item
  \textbf{Python-based Logic} --- The core engine (\texttt{humOS}) is
  implemented in Python and runs directly in the browser via Pyodide.
\end{itemize}

\subsection{Project Structure}\label{project-structure}

\begin{itemize}
\tightlist
\item
  \texttt{humos.py}: The core Universal Information Geometry Library.
\item
  \texttt{cell-src.js}: Source for the interactive code execution
  environment (Python REPL).
\item
  \texttt{cell.js}: Bundled and minified execution engine.
\item
  \texttt{*.html}: TDD documentation and interactive playgrounds.
\end{itemize}

\subsection{Interactive Demos}\label{interactive-demos}

The documentation includes interactive Python cells where you can run
\texttt{humOS} code directly.

\begin{itemize}
\tightlist
\item
  \textbf{PRISM}: Content addressing and similarity.
\item
  \textbf{Field Math}: Manifold operations and clustering.
\item
  \textbf{Graph Analysis}: Articulation points, pathologies, and health
  metrics.
\item
  \textbf{HUFD}: Holonomic Unified Field Dynamics and coherence
  tracking.
\end{itemize}

\subsection{Local Development}\label{local-development}

To run the site locally:

\begin{Shaded}
\begin{Highlighting}[]
\ExtensionTok{python} \AttributeTok{{-}m}\NormalTok{ http.server 8080}
\end{Highlighting}
\end{Shaded}

Then visit \texttt{http://localhost:8080}. Hosted on GitHub Pages from
\texttt{master}. The \texttt{.nojekyll} file ensures Python and other
assets are served correctly.

\end{document}
